\documentclass[10pt,letterpaper]{article}
\usepackage{cornell-notes}

\title{Chapter 13: Digital Meets Analog}
\author{}
\date{}
\setDocTitle{Chapter 13: Digital Meets Analog}
\setDocSubtitle{Electronics}
\setDocOwner{Electronics Cornell Notes}
\setDocDate{\today}
\setCornellCollection{Electronics Cornell Notes}
\setCornellUnitType{Chapter}
\setCornellUnitNumber{13}
\setCornellUnitTitle{Digital Meets Analog}
\hypersetup{pdftitle={Chapter 13: Digital Meets Analog},pdfsubject={Cornell notes},pdfauthor={}}
\begin{document}
\maketitle

\begin{CornellOverviewBox}{Concept}
Data conversion bridges continuous signals and numerical codes. DAC and ADC architectures trade resolution, rate, latency, noise, linearity, power, and reference demands. Sampling introduces aliasing; quantization and converter nonidealities create error; drivers and references determine whether data-sheet performance survives. PLLs and pseudorandom sequences extend the chapter into mixed-signal timing and test generation.
\end{CornellOverviewBox}
\begin{CornellOverviewBox}{Concept}
\textbf{Prerequisites.} Binary codes, op-amp precision/noise, filters and sampling concepts, references, logic timing, and basic probability.

\textbf{After studying this chapter, you should be able to:}
\begin{itemize}
  \item Distinguish resolution, accuracy, offset, gain, DNL, INL, monotonicity, ENOB, SNR, SFDR, and settling.
  \item Compare string, R--2R, current-steering, PWM, flash, SAR, integrating, and delta--sigma conversion.
  \item Design anti-alias, reference, driver, grounding, and clock paths for a converter system.
  \item Explain PLL blocks and LFSR-based pseudorandom/noise generation.
\end{itemize}\textbf{Central design questions.}
\begin{itemize}
  \item What signal bandwidth, amplitude, source impedance, latency, and accuracy must be preserved?
  \item How do reference, clock jitter, driver settling, and layout limit the converter?
  \item Which errors are static code errors, spectral errors, random noise, or alias products?
\end{itemize}\end{CornellOverviewBox}
\section{Cornell Notes}
\begin{CornellNotesTable}
\CornellNoteRow{13.1: Converter vocabulary}{Resolution is code granularity, not guaranteed accuracy. Offset shifts all codes; gain error changes slope; DNL measures step-size deviation; INL measures transfer deviation; missing codes and monotonicity are related but distinct. Dynamic measures include SNR, SINAD, ENOB, THD, SFDR, aperture jitter, and settling. State code format, reference convention, full-scale definition, sample rate, input frequency, and bandwidth.}
\CornellNoteRow{13.2: DAC architectures}{A resistor string is monotonic but scales in elements; R--2R ladders reuse two resistor values; current steering supports speed but needs matched sources and controlled switching; multiplying DACs operate with variable reference under limits. Voltage output requires an amplifier with settling and stability. Delta--sigma DACs shape quantization noise; PWM relies on time averaging and filtering. Choose by speed, glitch, monotonicity, output type, reference, and load.}
\CornellNoteRow{13.3: DAC applications}{Sources and coil/current drivers add output amplifier, reference, compliance, switching glitch, load dynamics, protection, and calibration. Multi-channel systems need crosstalk and update synchronization analysis. A precision code does not guarantee a precision output if reference, resistor, amplifier, ground, or load errors dominate. Define behavior during reset and code transition.}
\CornellNoteRow{13.4: Linearity}{Endpoint and best-fit INL use different reference lines. DNL below $-1$ LSB implies missing code; monotonicity requires nonnegative steps. Code-dependent glitch is a dynamic error not captured by dc INL. Histogram and slow-ramp tests expose code widths; spectral tests expose harmonic/spur behavior. Calibrate only when the error is stable and observable.}
\CornellNoteRow{13.5: Sampling}{A sampler produces spectral images at multiples of sample rate; any input content above Nyquist can alias into baseband. Anti-alias filtering must attenuate signals/noise at every folding band. Quantization step is approximately full-scale range divided by $2^N$. Ideal quantization SNR for a full-scale sine is about $6.02N+1.76\,\mathrm{dB}$, but real noise, distortion, clock jitter, and input drive reduce it.}
\CornellNoteRow{13.6: Flash ADCs}{Flash conversion uses many comparators in parallel for very low latency, with exponential comparator count, input capacitance, power, offset matching, and bubble-code challenges. Folding/interpolation reduce resource count. High-speed ADCs need controlled-impedance differential drive, common-mode control, low-jitter clocks, clean references, and careful layout.}
\CornellNoteRow{13.7: SAR ADCs}{A SAR converter uses a DAC and comparator to perform a binary search, typically one bit per decision. It offers moderate/high speed, low latency, and good power. The input often presents dynamic capacitive kickback during acquisition; the driver and RC filter must settle to required accuracy within acquisition time. Source resistance limits are frequency and resolution dependent.}
\CornellNoteRow{13.8: Integrating ADCs}{Single-, dual-, and multislope converters integrate input and reference over controlled intervals. Dual-slope conversion rejects component scale errors and periodic interference when timing is chosen well, at the cost of speed. Switch leakage, dielectric absorption, reference stability, integrator linearity, auto-zero, and clock timing determine accuracy. Voltage-to-frequency conversion similarly trades amplitude for pulse count.}
\CornellNoteRow{13.9: Delta--sigma}{Oversampling, a feedback modulator, noise shaping, and digital decimation move quantization noise out of band and average it. Order and oversampling ratio trade resolution, stability, bandwidth, and latency. Idle tones, digital-filter settling, multiplexing limitations, clocking, and input/reference drive matter. The output rate is not identical to independent-settled measurement rate.}
\CornellNoteRow{13.10: Architecture choice}{Flash prioritizes latency/rate, SAR balanced rate/power/latency, integrating precision and rejection, and delta--sigma high in-band resolution with latency. Compare real ENOB versus input frequency, bandwidth, channel switching, power modes, reference, clock, interface, calibration, and availability. Averaging reduces uncorrelated noise but not deterministic INL or coherent interference.}
\CornellNoteRow{13.11--13.12: Converter systems}{Integrated metering, touch, bridge, sequencer, and acquisition devices combine converters with muxes, PGAs, references, excitation, and digital logic. A multichannel system must specify per-channel settling, crosstalk, simultaneous versus sequential sampling, throughput, input protection, sensor bias, reference loading, and timestamp alignment. Partition analog ground by current flow, not by arbitrary cuts.}
\CornellNoteRow{13.13: Phase-locked loops}{A PLL includes phase/frequency detector, charge pump or detector gain, loop filter, VCO/NCO, and dividers. Loop dynamics set capture/lock, reference spur, jitter transfer, and VCO noise rejection. Multiplication raises output frequency and changes phase-noise relationships. Design with a linear model near lock, then verify nonlinear acquisition, divider states, and supply coupling.}
\CornellNoteRow{13.14: Pseudorandom sequences}{Maximal LFSRs generate deterministic sequences with long period and useful near-white discrete spectra after appropriate mapping/filtering. They support test, scrambling, dithering, and system identification but are not true entropy. Tap/polynomial convention, seed, period, balance, run lengths, clock images, and reconstruction filtering define the output. True random generation requires a physical entropy source and health testing.}
\end{CornellNotesTable}
\begin{CornellSummaryBox}{Summary}
Mixed-signal performance is a system result: converter architecture, reference, clock, driver, filter, layout, code handling, and timing all contribute. Resolution is only code count. Prevent aliasing before sampling, allow acquisition/settling, distinguish static and dynamic errors, and verify performance at the actual input frequency, amplitude, rate, and environment.
\end{CornellSummaryBox}
\section{Worked Examples}
\begin{CornellExamBox}{LSB and ideal SNR}
A 14-bit ADC spans $4.096\,\mathrm{V}$. One ideal LSB is $4.096/2^{14}=250\,\mu\mathrm{V}$. Ideal full-scale-sine quantization SNR is $6.02(14)+1.76=86.0\,\mathrm{dB}$. Neither number includes offset, reference, noise, INL, driver, or clock errors.
\end{CornellExamBox}
\begin{CornellExamBox}{Aperture-jitter limit}
For a $20\,\mathrm{MHz}$ sine and rms clock jitter $1.0\,\mathrm{ps}$, jitter-limited SNR is $-20\log_{10}(2\pi f\sigma_t)=78.0\,\mathrm{dB}$. To achieve 12 effective bits at that input, clock and aperture jitter likely need improvement plus margin for all other noise/distortion.
\end{CornellExamBox}
\begin{CornellExamBox}{SAR settling}
A $2\,\mathrm{V}$ step must settle within half an LSB of a 12-bit range, fractional error $1/8192$. A first-order driver requires about $-\ln(1/8192)=9.01$ time constants during acquisition. If acquisition is $2\,\mu\mathrm{s}$, $\tau$ should be no more than $222\,\mathrm{ns}$, before nonlinear recovery and ADC kickback.
\end{CornellExamBox}
\section{Design and Troubleshooting}
\begin{CornellExamBox}{Design and Troubleshooting}
\begin{enumerate}
  \item Specify signal range/bandwidth, source, channel count, sample timing, latency, accuracy, noise, and spectral requirements.
  \item Choose converter and reference/clock architecture; derive anti-alias, acquisition, settling, jitter, and data-throughput budgets.
  \item Design driver, differential/common-mode network, reference buffer/filter, grounding, decoupling, protection, and digital interface.
  \item Test dc transfer/histogram, noise, FFT, frequency/amplitude sweep, channel switching, reference/clock sensitivity, temperature, and data integrity.
\end{enumerate}\end{CornellExamBox}
\begin{CornellExamBox}{Symptoms, likely causes, and checks}
\begin{itemize}
  \item Missing/noisy codes: driver not settled, reference noise, grounding, input overrange, clock coupling, or digital read timing.
  \item Spurs move with input: distortion/aliasing; fixed spurs suggest clock, supply, interface, or coherent interference.
  \item Good single channel, poor scanning: mux charge memory, inadequate acquisition, crosstalk, or per-channel source impedance.
  \item DAC output glitches/oscillates: code-dependent switching, output amplifier recovery/stability, or reference disturbance.
  \item Delta--sigma appears stale: digital-filter latency, settling after mux/gain change, or misunderstood output-data rate.
\end{itemize}\end{CornellExamBox}
\begin{CornellWarningBox}{Safety and Limitations}
Converter inputs connected to field wiring require current-limited overvoltage and ESD protection that remains safe when unpowered. Isolation may be needed for high common-mode or hazardous systems. Never allow an ADC's small input range to encourage direct, unprotected connection to mains, shunts, batteries, or high-energy sensors.
\end{CornellWarningBox}
\section{Key-Equation Sheet}
\begin{CornellExamBox}{Key Equations}
\begin{itemize}
  \item Ideal LSB: $q=V_{FS}/2^N$ for the stated full-scale convention.
  \item Ideal sine quantization SNR: $6.02N+1.76\,\mathrm{dB}$.
  \item ENOB from SINAD: $(SINAD-1.76)/6.02$ bits.
  \item Jitter-limited sine SNR: $-20\log_{10}(2\pi f_{in}\sigma_t)$.
  \item First-order settling to fractional error $\epsilon$: $t=-\tau\ln\epsilon$.
  \item Nyquist condition prevents alias only when the input is sufficiently band-limited below $f_s/2$.
\end{itemize}\end{CornellExamBox}
\section{Study and Review}
\subsection{Design checklist}
\begin{itemize}
  \item Code format, full-scale convention, reference, and metric definitions are explicit.
  \item Anti-alias attenuation and acquisition/settling budgets pass every channel.
  \item Clock jitter and reference noise support target SNR/ENOB.
  \item Driver common mode, output current, RC network, and kickback are validated.
  \item Layout, decoupling, data timing, reset, calibration, and spectral test setup are controlled.
\end{itemize}\subsection{Common misconceptions}
\begin{itemize}
  \item Nominal bits do not equal ENOB or accuracy.
  \item Sampling faster does not remove the need for anti-alias filtering.
  \item Averaging does not repair INL or coherent spurs.
  \item The ADC input is not always a static high resistance.
  \item Delta--sigma output rate does not equal zero-latency independent samples.
\end{itemize}\subsection{Active-recall questions}
\begin{enumerate}
  \item Define DNL and INL.
  \item What makes a DAC monotonic?
  \item Why does a SAR input kick back charge?
  \item How does dual-slope reject component error?
  \item What is noise shaping?
  \item How does clock jitter limit high-frequency SNR?
  \item What determines anti-alias filter attenuation?
  \item How do flash and SAR latency differ?
  \item What does a PLL loop bandwidth trade?
  \item Why is an LFSR not true random?
\end{enumerate}\subsection{Original practice problems}
\begin{enumerate}
  \item Find ideal LSB for a $2.5\,\mathrm{V}$, 16-bit range.
  \item An ADC has $SINAD=68\,\mathrm{dB}$. Estimate ENOB.
  \item Find jitter-limited SNR for $f=5\,\mathrm{MHz}$ and $\sigma_t=2\,\mathrm{ps}$.
\end{enumerate}\subsection{Answer key / solution outline}
\begin{enumerate}
  \item $2.5/65536=38.15\,\mu\mathrm{V}$.
  \item $(68-1.76)/6.02=11.0$ bits.
  \item $-20\log_{10}(2\pi(5\,\mathrm{MHz})(2\,\mathrm{ps}))=84.0\,\mathrm{dB}$.
\end{enumerate}\begin{CornellSummaryBox}{Summary}
Bits are not performance. Preserve the analog signal with anti-aliasing, acquisition settling, quiet reference and clock, correct common mode, and disciplined layout; then measure static transfer and dynamic spectrum under real operating conditions.
\end{CornellSummaryBox}

\end{document}
